\frfilename{mt2a5.tex}
\versiondate{13.11.07}
\copyrightdate{1996}

\def\chaptername{Lampiran}
\def\sectionname{Ruang topologis linear}

\newsection{2A5}

Tujuan utama \S2A3 sesungguhnya adalah mengkaji topologi-topologi
tertentu pada ruang-ruang linear dalam Bab 24. Di sini saya memberikan
beberapa bagian dari teori umumnya.

\leader{2A5A}{Topologi \dvrocolon{ruang linear}}\cmmnt{ Salah satu hal
yang tidak dibahas secara terperinci dalam setiap pengantar analisis
fungsional adalah konsep umum ruang topologis linear. Gagasan-gagasan
yang diperlukan untuk pembahasan \S245 dapat dinyatakan dengan cukup
ringkas.

\medskip

\noindent}{\bf Definisi} Sebuah {\bf ruang topologis linear} atau
{\bf ruang vektor topologis} atas $\RoverC$ adalah ruang linear $U$
atas $\RoverC$ yang dilengkapi dengan topologi $\frak T$ sedemikian
sehingga kedua peta

\Centerline{$(u,v)\mapsto u+v:U\times U\to U$,}

\Centerline{$(\alpha,u)\mapsto \alpha u:\RoverC\times U\to U$}

\noindent kontinu, dengan ruang produk $U\times U$ dan
$\RoverC\times U$ diberi topologi produknya\cmmnt{ (2A3T)}. Diberikan
ruang linear $U$, sebuah topologi pada $U$ yang memenuhi syarat-syarat
di atas disebut {\bf topologi ruang linear}. Perhatikan bahwa

\Centerline{$(u,v)\mapsto u-v=u+(-1)v:U\times U\to U$}

\noindent juga akan kontinu.

\leader{2A5B}{}\cmmnt{ Semua ruang topologis linear yang kita perlukan
ternyata dapat disajikan dengan mudah melalui rumusan berikut.

\medskip

\noindent}{\bf Proposisi} Misalkan $U$ adalah ruang linear atas
$\RoverC$, dan $\Tau$ adalah keluarga fungsional
$\tau:U\to\coint{0,\infty}$ sedemikian sehingga

\quad (i) $\tau(u+v)\le\tau(u)+\tau(v)$ untuk semua $u$, $v\in U$,
$\tau\in\Tau$;

\quad (ii) $\tau(\alpha u)\le\tau(u)$ jika $u\in U$, $|\alpha|\le 1$,
$\tau\in\Tau$;

\quad (iii) $\lim_{\alpha\to 0}\tau(\alpha u)=0$ untuk setiap $u\in U$,
$\tau\in \Tau$.

\noindent Untuk $\tau\in\Tau$, definisikan
$\rho_{\tau}:U\times U\to\coint{0,\infty}$ dengan menetapkan
$\rho_{\tau}(u,v)=\tau(u-v)$ untuk semua $u$, $v\in U$. Maka setiap
$\rho_{\tau}$ adalah pseudometrik pada $U$, dan topologi yang
didefinisikan oleh $\Rho=\{\rho_{\tau}:\tau\in\Tau\}$ menjadikan $U$
ruang topologis linear.

\proof{{\bf (a)} Patut segera dicatat bahwa

\Centerline{$\tau(0)=\lim_{\alpha\to 0}\tau(\alpha 0)=0$}

\noindent untuk setiap $\tau\in\Tau$.

\medskip

{\bf (b)} Untuk melihat bahwa setiap $\rho_{\tau}$ adalah
pseudometrik, gunakan argumen berikut.

\medskip

\quad{\bf (i)} $\rho_{\tau}$ mengambil nilai dalam
$\coint{0,\infty}$ karena $\tau$ juga demikian.

\medskip

\quad{\bf (ii)} Jika $u$, $v$, $w\in U$, maka

$$\eqalign{\rho_{\tau}(u,w)
&=\tau(u-w)
=\tau((u-v)+(v-w))\cr
&\le\tau(u-v)+\tau(v-w)
=\rho_{\tau}(u,v)+\rho_{\tau}(v,w).\cr}$$

\medskip

\quad{\bf (iii)} Jika $u$, $v\in U$, maka

\Centerline{$\rho_{\tau}(v,u)=\tau(v-u)
=\tau(-1(u-v))\le\tau(u-v)=\rho_{\tau}(u,v)$,}

\noindent dan secara serupa
$\rho_{\tau}(u,v)\le\rho_{\tau}(v,u)$, sehingga keduanya sama.

\medskip

\quad{\bf (iv)} Jika $u\in U$, maka
$\rho_{\tau}(u,u)=\tau(0)=0$.

\medskip

{\bf (c)} Misalkan $\frak T$ adalah topologi pada $U$ yang
didefinisikan oleh $\{\rho_{\tau}:\tau\in\Tau\}$ (2A3F).

\medskip

\quad{\bf (i)} Penjumlahan kontinu karena, diberikan $\tau\in\Tau$,
kita mempunyai

$$\eqalign{\rho_{\tau}(u'+v',u+v)
&=\tau((u'+v')-(u+v))\cr
&\le\tau(u'-u)+\tau(v'-v)
\le\rho_{\tau}(u',u)+\rho_{\tau}(v',v)\cr}$$

\noindent untuk semua $u$, $v$, $u'$, $v'\in U$. Ini berarti bahwa,
diberikan $\epsilon>0$ dan $(u,v)\in U\times U$, kita akan mempunyai

\Centerline{$\rho_{\tau}(u'+v',u+v)\le\epsilon$ apabila
$(u',v')\in
U((u,v);\tilde\rho_{\tau},\bar\rho_{\tau};\Bover{\epsilon}2)$,}

\noindent dengan menggunakan bahasa 2A3Tb. Karena
$\tilde\rho_{\tau}$, $\bar\rho_{\tau}$ adalah dua dari pseudometrik
yang mendefinisikan topologi produk $U\times U$ (2A3Tb), maka
$(u,v)\mapsto u+v$ kontinu menurut kriteria 2A3H.

\medskip

\quad{\bf (ii)} Perkalian skalar kontinu karena jika $u\in U$ dan
$n\in\Bbb N$, maka $\tau(nu)\le n\tau(u)$ untuk setiap
$\tau\in\Tau$ (gunakan induksi pada $n$). Akibatnya, jika
$\tau\in\Tau$,

\Centerline{$\tau(\alpha u)\le n\tau(\Bover{\alpha}{n}u)\le n\tau(u)$}

\noindent apabila $|\alpha|<n\in\Bbb N$ dan $\tau\in\Tau$. Sekarang,
diberikan $(\alpha,u)\in \RoverC\times U$ dan $\epsilon>0$, ambil
$n>|\alpha|$ dan $\delta>0$ sedemikian sehingga
$\delta\le\min(n-|\alpha|,\bover{\epsilon}{2n})$ dan
$\tau(\gamma u)\le\bover{\epsilon}2$ apabila
$|\gamma|\le\delta$; maka

$$\eqalign{\rho_{\tau}(\alpha'u',\alpha u)
&=\tau(\alpha'u'-\alpha u)
\le\tau(\alpha'(u'-u))+\tau((\alpha'-\alpha)u)\cr
&\le n\tau(u'-u)+\tau((\alpha'-\alpha)u)\cr}$$

\noindent apabila $u'\in U$, $\alpha'\in\RoverC$, dan
$|\alpha'|<n\in\Bbb N$. Karena itu, dengan menetapkan
$\theta(\alpha',\alpha)=|\alpha'-\alpha|$ untuk $\alpha'$,
$\alpha\in\RoverC$,

\Centerline{$\rho_{\tau}(\alpha'u',\alpha u)
\le n\delta+\Bover{\epsilon}2\le\epsilon$}

\noindent apabila

\Centerline{$(\alpha',u')\in
U((\alpha,u);\tilde\theta,\bar\rho_{\tau};\delta)$.}

\noindent Karena $\tilde\theta$ dan $\bar\rho_{\tau}$ termasuk di
antara pseudometrik yang mendefinisikan topologi
$\RoverC\times U$, peta $(\alpha,u)\mapsto\alpha u$ memenuhi kriteria
2A3H dan kontinu.

Jadi $\frak T$ adalah topologi ruang linear pada $U$.
}%end of proof of  2A5B

\medskip

\noindent{\bf Catatan}\dvAnew{2013}
Fungsional yang memenuhi syarat (i)--(iii) di atas disebut
{\bf F-seminorma}; sebuah F-seminorma $\tau$ sedemikian sehingga
$\tau(u)\ne 0$ untuk setiap $u$ tak nol disebut {\bf F-norma}.

\leader{*2A5C}{}\cmmnt{ Kita tidak memerlukannya untuk Bab 24, tetapi
hasil berikut patut diketahui.

\medskip

\noindent}{\bf Teorema} Misalkan $U$ adalah ruang linear dan
$\frak T$ topologi ruang linear pada $U$.

(a) Terdapat keluarga F-seminorma $\Tau$ yang mendefinisikan
$\frak T$ seperti dalam 2A5B.

(b) Jika $\frak T$ dapat dimetrikkan, kita dapat mengambil $\Tau$
yang hanya terdiri atas satu fungsional.

\proof{{\bf (a)} {\smc Kelley \& Namioka 76}, hlm.\ 50.

\medskip

{\bf (b)} {\smc K\"othe 69}, \S15.11.
}%end of proof of 2A5C

\vleader{72pt}{2A5D}{Definisi} Misalkan $U$ adalah ruang linear atas
$\RoverC$. Sebuah {\bf seminorma} pada $U$ adalah fungsional
$\tau:U\to\coint{0,\infty}$ sedemikian sehingga

\quad (i) $\tau(u+v)\le\tau(u)+\tau(v)$ untuk semua $u$, $v\in U$;

\quad (ii) $\tau(\alpha u)=|\alpha|\tau(u)$ jika $u\in U$,
$\alpha\in\RoverC$.

\cmmnt{\noindent Perhatikan bahwa norma selalu merupakan seminorma,
dan seminorma selalu merupakan F-seminorma. Khususnya, pengaitan suatu
metrik dengan norma (2A4Bb) merupakan kasus khusus dari 2A5B.}

\leader{2A5E}{Himpunan konveks (a)} Misalkan $U$ adalah ruang linear
atas $\RoverC$. Suatu himpunan bagian $C$ dari $U$ disebut
{\bf konveks} jika $\alpha u+(1-\alpha)v\in C$ apabila $u$, $v\in C$
dan $\alpha\in[0,1]$. Irisan sebarang keluarga himpunan konveks adalah
konveks, sehingga untuk setiap himpunan $A\subseteq U$ terdapat
himpunan konveks terkecil yang memuat $A$; himpunan ini\cmmnt{ tepat}
merupakan himpunan vektor yang dapat dinyatakan sebagai
$\sum_{i=0}^n\alpha_iu_i$, dengan
$u_0,\ldots,u_n\in A$, $\alpha_0,\ldots,\alpha_n\in[0,1]$, dan
$\sum_{i=0}^n\alpha_i=1$\prooflet{ ({\smc Bourbaki 87}, II.2.3)};
himpunan ini disebut {\bf selubung konveks} dari $A$. Jika
$C$, $C'\subseteq U$ konveks, dan $\alpha\in\RoverC$, maka
$\alpha C$ dan $C+C'$ konveks\dvAnew{2012}.

\spheader 2A5Eb Jika $U$ adalah ruang topologis linear, penutupan
sebarang himpunan konveks adalah konveks\prooflet{
({\smc Bourbaki 87}, II.2.6)}. Akibatnya, untuk setiap $A\subseteq U$,
penutupan selubung konveks dari $A$ merupakan himpunan konveks tertutup
terkecil yang memuat $A$; himpunan ini disebut
{\bf selubung konveks tertutup} dari $A$.

\spheader 2A5Ec Untuk rujukan mendatang, saya mencatat bahwa dalam
ruang topologis linear, penutupan sebarang subruang linear adalah
subruang linear.
\prooflet{({\smc Bourbaki 87}, I.1.3; {\smc K\"othe 69}, \S15.2.
Bandingkan 2A4Cb.)}

\leader{2A5F}{Kelengkapan dalam ruang topologis
\dvrocolon{linear}}\cmmnt{ Dalam ruang bernorma, kelengkapan dapat
diuraikan melalui barisan Cauchy (2A4D). Dalam ruang topologis linear
umum, hal ini tidak memadai. Teori kelengkapan yang sebenarnya
memerlukan konsep ruang seragam (lihat \S3A4 dalam jilid berikutnya,
atau {\smc Kelley 55}, bab 6; {\smc Engelking 89}, \S8.1;
{\smc Bourbaki 66}, bab II); saya tidak akan menguraikannya di sini,
tetapi akan memberikan versi yang disesuaikan dengan ruang linear.
Saya menyebutkan hal ini hanya karena saya berharap suatu hari Anda
akan sampai pada teori umum (dalam Jilid 3 risalah ini, jika tidak lebih
awal), dan Anda perlu menyadari bahwa kasus khusus yang diuraikan di
sini memberikan penekanan yang menyesatkan pada beberapa bagian.

\medskip

\noindent}{\bf Definisi} Misalkan $U$ adalah ruang linear atas
$\RoverC$, dan $\frak T$ topologi ruang linear pada $U$. Sebuah filter
$\Cal F$ pada $U$ disebut {\bf Cauchy} jika untuk setiap himpunan
terbuka $G$ di $U$ yang memuat $0$, terdapat $F\in\Cal F$ sedemikian
sehingga $F-F=\{u-v:u,\,v\in F\}$ termuat dalam $G$.
$U$ disebut {\bf lengkap} jika setiap filter Cauchy pada $U$ konvergen.

\leader{2A5G}{}\cmmnt{ Filter Cauchy mempunyai uraian sederhana apabila
topologi ruang linear didefinisikan dengan metode 2A5B.

\medskip

\noindent}{\bf Lema} Misalkan $U$ adalah ruang linear atas $\RoverC$,
dan $\Tau$ keluarga F-seminorma yang mendefinisikan topologi ruang
linear pada $U$, seperti dalam 2A5B. Maka sebuah filter $\Cal F$ pada
$U$ adalah Cauchy jika dan hanya jika untuk setiap $\tau\in\Tau$ dan
$\epsilon>0$ terdapat $F\in\Cal F$ sedemikian sehingga
$\tau(u-v)\le\epsilon$ untuk semua $u$, $v\in F$.

\proof{{\bf (a)} Misalkan $\Cal F$ Cauchy, $\tau\in\Tau$, dan
$\epsilon>0$. Maka $G=U(0;\rho_{\tau};\epsilon)$ terbuka (dengan
menggunakan bahasa 2A3F--2A3G), sehingga terdapat $F\in\Cal F$
sedemikian sehingga $F-F\subseteq G$; tetapi ini tepat berarti bahwa
$\tau(u-v)<\epsilon$ untuk semua $u$, $v\in F$.

\medskip

{\bf (b)} Misalkan $\Cal F$ memenuhi kriteria tersebut, dan $G$
merupakan himpunan terbuka yang memuat $0$. Maka terdapat
$\tau_0,\ldots,\tau_n\in\Tau$ dan $\epsilon>0$ sedemikian sehingga
$U(0;\rho_{\tau_0},\ldots,\rho_{\tau_n};\epsilon)\subseteq G$.
Untuk setiap $i\le n$, terdapat $F_i\in\Cal F$ sedemikian sehingga
$\tau_i(u-v)<\bover{\epsilon}2$ untuk semua $u$, $v\in F_i$; sekarang
$F=\bigcap_{i\le n}F_i\in\Cal F$ dan $u-v\in G$ untuk semua
$u$, $v\in F$.
}%end of proof of 2A5G

\leader{2A5H}{Ruang bernorma dan kelengkapan
\dvrocolon{sekuensial}}\cmmnt{ Perlu saya tunjukkan bahwa untuk ruang
bernorma, definisi 2A5F sesuai dengan definisi 2A4D.

\medskip

\noindent}{\bf Proposisi} Misalkan $(U,\|\,\|)$ adalah ruang bernorma
atas $\RoverC$, dan $\frak T$ topologi ruang linear pada $U$ yang
didefinisikan dengan metode 2A5B dari himpunan
$\Tau=\{\|\,\|\}$. Maka $U$ lengkap dalam pengertian 2A5F jika dan
hanya jika lengkap dalam pengertian 2A4D.

\proof{{\bf (a)} Pertama-tama, misalkan $U$ lengkap dalam pengertian
2A5F. Misalkan $\sequencen{u_n}$ adalah barisan di $U$ yang Cauchy
dalam pengertian 2A4Da. Tetapkan

\Centerline{$\Cal F
=\{F:F\subseteq U,\,\{n:u_n\notin F\}\text{ berhingga}\}$.}

\noindent Maka mudah diperiksa bahwa $\Cal F$ adalah filter pada $U$,
yaitu citra filter Fr\'echet di bawah peta
$n\mapsto u_n:\Bbb N\to U$. Jika $\epsilon>0$, ambil $m\in\Bbb N$
sedemikian sehingga $\|u_j-u_k\|\le\epsilon$ apabila $j$, $k\ge m$;
maka $F=\{u_j:j\ge m\}$ termasuk dalam $\Cal F$, dan
$\|u-v\|\le\epsilon$ untuk semua $u$, $v\in F$. Jadi $\Cal F$ Cauchy
dalam pengertian 2A5F, dan mempunyai suatu limit, katakanlah $u$.
Sekarang, untuk setiap $\epsilon>0$, himpunan
$\{v:\|v-u\|<\epsilon\}=U(u;\rho_{\|\,\|};\epsilon)$ merupakan
himpunan terbuka yang memuat $u$, sehingga termasuk dalam $\Cal F$,
dan $\{n:\|u_n-u\|\ge\epsilon\}$ berhingga, yaitu, terdapat
$m\in\Bbb N$ sedemikian sehingga $\|u_n-u\|<\epsilon$ apabila
$n\ge m$. Karena $\epsilon$ sebarang,
$u=\lim_{n\to\infty}u_n$ dalam pengertian 2A3M. Karena
$\sequencen{u_n}$ sebarang, $U$ lengkap dalam pengertian 2A4D.

\medskip

{\bf (b)} Sekarang misalkan $U$ lengkap dalam pengertian 2A4D.
Misalkan $\Cal F$ adalah filter Cauchy pada $U$. Untuk setiap
$n\in\Bbb N$, pilih himpunan $F_n\in\Cal F$ sedemikian sehingga
$\|u-v\|\le 2^{-n}$ untuk semua $u$, $v\in F_n$. Untuk setiap
$n\in\Bbb N$, $F'_n=\bigcap_{i\le n}F_i$ termasuk dalam $\Cal F$,
sehingga tidak kosong; pilih $u_n\in F'_n$. Jika $m\in\Bbb N$ dan
$j$, $k\ge m$, maka $u_j$ dan $u_k$ keduanya termasuk dalam $F_m$,
sehingga $\|u_j-u_k\|\le 2^{-m}$; jadi $\sequencen{u_n}$ merupakan
barisan Cauchy dalam pengertian 2A4Da, dan mempunyai suatu limit,
katakanlah $u$. Sekarang ambil sebarang $\epsilon>0$ dan
$m\in\Bbb N$ sedemikian sehingga $2^{-m+1}\le\epsilon$. Tentu terdapat
$k\ge m$ sedemikian sehingga $\|u_k-u\|\le 2^{-m}$; sekarang
$u_k\in F_m$, sehingga

\Centerline{$F_m\subseteq\{v:\|v-u_k\|\le 2^{-m}\}
\subseteq\{v:\|v-u\|\le 2^{-m+1}\}
\subseteq\{v:\rho_{\|\,\|}(v,u)\le\epsilon\}$,}

\noindent dan
$\{v:\rho_{\|\,\|}(v,u)\le\epsilon\}\in\Cal F$.
Karena $\epsilon$ sebarang, $\Cal F$ konvergen ke $u$, menurut 2A3Sd.
Karena $\Cal F$ sebarang, $U$ lengkap.

\medskip

{\bf (c)}
Jadi kedua definisi itu bertepatan, setidaknya dengan syarat kita
memperbolehkan banyaknya pilihan serentak yang terhitung untuk $u_n$
dalam bagian (b) pembuktian.
}%end of proof of 2A5H

\leader{2A5I}{Topologi lemah}\cmmnt{ Sekarang saya beralih ke catatan
ringkas tentang topologi lemah pada ruang bernorma; dari sudut pandang
jilid ini, topologi-topologi tersebut sesungguhnya merupakan contoh
utama topologi ruang linear.}
Misalkan $U$ adalah ruang linear bernorma atas $\RoverC$.

\spheader 2A5Ia Tuliskan $U^*$ untuk dualnya
$\eurm B(U;\RoverC)$\cmmnt{ (2A4H)}. \cmmnt{Jika $h\in U^*$, maka
$|h|:U\to\coint{0,\infty}$ adalah seminorma, sehingga}
$\Tau=\{|h|:h\in U^*\}$ mendefinisikan topologi ruang linear pada
$U$\cmmnt{, menurut 2A5B}; topologi ini\cmmnt{ disebut}
{\bf topologi lemah} dari $U$.

\spheader 2A5Ib Sebuah filter $\Cal F$ pada $U$ konvergen ke $u\in U$
untuk topologi lemah dari $U$ jika dan hanya jika\cmmnt{
$\lim_{v\to\Cal F}\rho_{|h|}(v,u)=0$ untuk setiap $h\in U^*$ (2A3Sd),
yaitu, jika dan hanya jika $\lim_{v\to\Cal F}|h(v-u)|=0$ untuk setiap
$h\in U^*$, yaitu, jika dan hanya jika}
$\lim_{v\to\Cal F}h(v)=h(u)$ untuk setiap $h\in U^*$.

\spheader 2A5Ic Sebuah himpunan $C\subseteq U$ disebut
{\bf kompak lemah} jika kompak untuk topologi lemah dari $U$. Jadi
\cmmnt{ (dengan tunduk pada aksioma pilihan)} suatu himpunan
$C\subseteq U$ kompak lemah jika dan hanya jika untuk setiap ultrafilter
$\Cal F$ pada $U$ yang memuat $C$, terdapat $u\in C$ sedemikian
sehingga $\lim_{v\to\Cal F}h(v)=h(u)$ untuk setiap
$h\in U^*$\cmmnt{ (gabungkan 2A3R dengan (b) di atas)}.

\spheader 2A5Id Suatu himpunan bagian $A$ dari $U$ disebut
{\bf relatif kompak lemah} jika merupakan himpunan bagian dari suatu
himpunan bagian $U$ yang kompak lemah.

\spheader 2A5Ie Jika $h\in U^*$, maka $h:U\to\RoverC$ kontinu untuk
topologi lemah pada $U$ dan topologi biasa pada $\RoverC$\cmmnt{;
hal ini langsung terlihat jika kita menerapkan kriteria 2A3H}. Jadi,
jika $A\subseteq U$ relatif kompak lemah, $h[A]$ harus terbatas dalam
$\RoverC$. \prooflet{\Prf\ Misalkan $C\supseteq A$ adalah himpunan
kompak lemah. Maka $h[C]$ kompak dalam $\RoverC$, menurut 2A3Nb, sehingga
terbatas, menurut 2A2F (dengan mencatat bahwa jika lapangan dasarnya
adalah $\Bbb C$, maka lapangan itu dapat diidentifikasi, sebagai ruang
metrik, dengan $\BbbR^2$). Karena itu $h[A]$ juga terbatas.\ \Qed}

\spheader 2A5If Jika $V$ adalah ruang bernorma lain dan $T:U\to V$
operator linear terbatas, maka $T$ kontinu untuk masing-masing topologi
lemah. \prooflet{\Prf\ Jika $h\in V^*$, maka komposisi $hT$ termasuk
dalam $U^*$. Sekarang, untuk sebarang $u$, $v\in U$,

\Centerline{$\rho_{|h|}(Tu,Tv)=|h(Tu-Tv)|=|hT(u-v)|=\rho_{|hT|}(u,v)$,}

\noindent dengan $\rho_{|h|}$, $\rho_{|hT|}$ masing-masing merupakan
pseudometrik pada $V$, $U$ yang didefinisikan oleh rumus 2A5B.
Menurut 2A3H, $T$ kontinu.\ \Qed}

\spheader 2A5Ig Bersesuaian dengan topologi lemah pada ruang bernorma
$U$, kita mempunyai topologi {\bf lemah*} atau {\bf w*-} pada dualnya
$U^*$, yang didefinisikan oleh himpunan
$\Tau=\{|\hat u|:u\in U\}$, dengan saya menuliskan
$\hat u(f)=f(u)$ untuk setiap $f\in U^*$, $u\in U$. Seperti dalam (a),
ini adalah topologi ruang linear pada $U^*$. \cmmnt{(Penting untuk
membedakan topologi lemah* dari topologi lemah pada $U^*$. Yang pertama
hanya bergantung pada aksi $U$ pada $U^*$, sedangkan yang kedua
bergantung pada aksi $U^{**}=(U^*)^*$. Anda tidak akan mengalami
kesulitan untuk memeriksa bahwa $\hat u\in U^{**}$ untuk setiap
$u\in U$, tetapi pokoknya adalah bahwa mungkin terdapat anggota
$U^{**}$ yang tidak dapat direpresentasikan dengan cara ini, sehingga
muncul himpunan terbuka untuk topologi lemah yang tidak terbuka untuk
topologi lemah*.)}

\leader{*2A5J}{Ruang angelik}\cmmnt{ Saya tidak bergantung pada
gagasan-gagasan berikut, tetapi gagasan ini dapat memperjelas beberapa
hasil dalam \S\S246-247.} Pertama, sebuah ruang topologis $X$ disebut
{\bf reguler} jika apabila $G\subseteq X$ terbuka dan $x\in G$, maka
terdapat himpunan terbuka $H$ sedemikian sehingga
$x\in H\subseteq\overline{H}\subseteq G$. Berikutnya, sebuah ruang
Hausdorff reguler $X$ disebut {\bf angelik} jika apabila $A\subseteq X$
sedemikian sehingga setiap barisan dalam $A$ mempunyai titik gugus di
$X$, maka $\overline{A}$ kompak dan setiap titik $\overline{A}$ adalah
limit dari suatu barisan dalam $A$.\cmmnt{ Artinya, kekompakan dalam
$X$, serta topologi himpunan bagian kompak dari $X$, dapat diuraikan
secara efektif melalui barisan.} Sekarang\cmmnt{ teorema (oleh Eberlein
dan \v Smulian) menyatakan bahwa} setiap ruang bernorma bersifat angelik
dalam topologi lemahnya.\cmmnt{ (462D dalam Jilid 4;
{\smc K\"othe 69}, \S24; {\smc Dunford \& Schwartz 57}, V.6.1.)
Khususnya, hal ini berlaku bagi ruang-ruang $L^1$, sehingga tidak terlalu
mengejutkan bahwa terdapat kriteria untuk kekompakan lemah dalam ruang
$L^1$ yang hanya membahas barisan.
}%end of comment

\discrpage
